\begin{prob}
    Determine the asymptotic time complexity of the following piece of code, showing your reasoning
    and your work.

    \inputminted{python}{\thisdir/include/code.py}

    \textbf{Hint}: you might need to think back to calculus to remember the formula for
    the sum of a geometric progression... or you can check wikipedia.%
    \footnote{\url{https://en.wikipedia.org/wiki/Geometric_series}}

    \begin{soln}
        $\Theta(n^2)$.

        How many times does the \python{print} statement execute? That will
        tell us the time complexity, because each line of code takes constant
        time to run once, and the \python{print} statement is the line that is
        run most.

        On the first iteration
        of the outer loop, \python{i = 2} and this line runs $i^2 = 2^2 = 4$ times. On the second iteration
        of the outer loop, \python{i = 4} and this line runs $4^2 = 16$ times. On the third
        iteration, it runs 64 times. The pattern is captured by the table below:

        % a table with columns:
        %   - iteration #
        %   - $i$
        %   - $i^2$
        \begin{center}
            \begin{tabular}{ccc}
                \toprule
                Iteration \# & $i$      & \# of \python{print}s \\
                \midrule
                1            & 2        & 4                     \\
                2            & 4        & 16                    \\
                3            & 8        & 64                    \\
                4            & 16       & 256                   \\
                $\vdots$     & $\vdots$ & $\vdots$              \\
                $k$          & $2^k$    & $2^{2k} = 4^k$        \\
                $\vdots$     & $\vdots$ & $\vdots$              \\
                \bottomrule
            \end{tabular}
        \end{center}

        The third column tells us the number of times \python{print} is
        executed \emph{during} that iteration of the \python{while} loop. In
        general, on the $k$th iteration of the outer loop, the \python{print}
        runs $2^{2k}$ times, which happens to equal $4^k$. The total number of
        \python{print}s will be the sum of the third column.

        We need to sum up the third column, but we don't yet know the last term
        in that sum. To figure it out, recognize that the outer loop will run
        \emph{roughly} $\log_2 n$ times since $i$ is doubling on each
        iteration. We'll see in a moment why it's OK to have a rough estimate
        instead of being exact here.

        The total number of executions of the \python{print} is therefore:
        \[
            \underbrace{4}_{\text{\footnotesize 1st outer iter.}}
            +
            \underbrace{16}_{\text{\footnotesize 2nd outer iter.}}
            +
            \underbrace{64}_{\text{\footnotesize 3rd outer iter.}}
            +
            \ldots
            +
            \underbrace{4^{k}}_{\text{\footnotesize $k$th outer iter.}}
            +
            \ldots
            +
            \underbrace{4^{\log_2 n}}_{\text{\footnotesize $\log_2 n$th outer iter.}}
        \]

        Or, written using summation notation:
        \[
            \sum_{k=1}^{\log_2 n} 4^k
        \]

        This is the sum of a \emph{geometric progression}. Wikipedia tells us that the formula
        for the sum of $1 + 4 + 8 + 16 + \ldots + 4^K$ is:
        \[
            \sum_{k=0}^{K} 4^k = \frac{1 - 4^{K+1}}{1 - 4} = \frac{4^{K+1} -1}{3}
        \]
        Notice that this sum starts from $k = 0$, while ours starts with $k = 1$.
        In other words, our sum is the same except it is missing the first term corresponding
        to $k = 0$.
        So to
        ``correct'' this, we subtract the $k = 0$ term (which is $4^0 = 1$)
        from the larger sum :
        \[
            \sum_{k=1}^{K} 4^k =
            \left(\sum_{k=0}^{K} 4^k\right) - 4^0 =
            \frac{4^{K+1} -1}{3} - 1 =
            \frac{4^{K+1} - 4}{3}
        \]
        Plugging in $K = \log_2 n$:
        \begin{align*}
            \sum_{k=1}^{\log_2 n} 4^k
             & = \frac{4^{\log_2 n + 1} - 4}{3}
            \intertext{Using the fact that $a^{b+c} = a^b \cdot a^c$:}
             & =
            \frac{4 \cdot 4^{\log_2 n} - 4}{3}           \\
             & =
            \frac{4 \cdot 4^{\log_2 n}}{3} - \frac{4}{3} \\
             & =
            \Theta(4^{\log_2 n})                         \\
            \intertext{Simplifying using the fact that $ (a^b)^c = a^{b \cdot c}$:}
             & =
            \Theta((2^2)^{\log_2 n})                     \\
             & =
            \Theta((2^{\log_2 n})^2)                     \\
             & =
            \Theta(n^2)
        \end{align*}

        Now, remember that we said that we could afford to be a little imprecise when calculating
        the number of times that the outer loop runs. Let's see why. If $n$ isn't a power
        of $2$, $\log_2 n$ will not be an integer. For instance, if $n = 17$, $\log_2 n = 4.08$.
        Of course, the loop can't iterate 4.08 times -- you can check that it will actually iterate
        5 times. In general, the loop will run exactly $\lceil \log_2 n \rceil$ times, where
        $\lceil \cdot \rceil$ is the \emph{ceiling} operation; it rounds a real number up
        to the next integer.

        In other words, $\log_2 n$ could actually be as much as one less than the actual number
        of iterations. Perhaps to be careful we should overestimate the number of iterations
        to be $\log_2 n + 1$. What if we were to use $\log_2 n + 1$ instead? We'd end up with:
        an extra term in the sum, which we can split off:
        \[
            \sum_{k=1}^{\log_2 n + 1} 4^k = \left(
            \sum_{k=1}^{\log_2 n} 4^k
            \right)
            +
            4^{\log_2 n + 1}
        \]
        The part in parentheses is the sum we calculated above, which we know is $\Theta(n^2)$.
        The ``extra'' term, $4^{\log_2 n + 1}$, is equal to $4 \cdot 4^{\log_2 n}$, which is
        $4 n^2$. So the total is still $$\Theta(n^2) + 4 n^2 = \Theta(n^2)$$
        and the time complexity doesn't change. This is why using $\Theta$ notation allows us
        to be ``sloppy'' at times without being incorrect. It saves us work, as long as
        we know how to use it correctly!
    \end{soln}

\end{prob}
